\documentclass[aspectratio=169,10pt,xcolor={dvipsnames}]{beamer}
\usepackage[utf8]{inputenc}
\usepackage[T1]{fontenc}
\usepackage{lmodern, textcomp}
\usepackage{amsmath, amssymb}
\usepackage{appendix}
\usepackage{babel}
\usepackage{ragged2e}
\usepackage{graphics}
\usepackage{dsfont}
\usepackage{subcaption}
\usepackage{chngcntr}
\usepackage{mathabx}
\usepackage{soul}
\usepackage{amsmath,amsthm}
\usepackage{multirow}
\usepackage{booktabs, threeparttable}
\usepackage{tabularx,booktabs}
\usepackage{ltablex,booktabs}
\usepackage{bm}
\counterwithin{subfigure}{figure}
\usetheme{default}
\usepackage{appendixnumberbeamer}
\usefonttheme[onlymath]{serif}
%\usecolortheme{spruce}

%Colours
\definecolor{shaded}{RGB}{153, 194, 173}
\definecolor{violet}{RGB}{51, 51, 178}
\definecolor{titlecolor}{RGB}{73, 156, 94}
\definecolor{titlecolor2}{RGB}{10, 74, 38}
\definecolor{alerted}{RGB}{51, 133, 92}
\definecolor{alerted1}{RGB}{229, 240, 235}
\definecolor{box}{RGB}{196, 224, 199}
\setbeamercolor{alerted text}{fg=alerted}
\setbeamercolor{section in toc}{fg=black}
\setbeamercolor{item}{fg=shaded}
\setbeamercolor{block body alerted}{use=structure,fg=black,bg=alerted1}
\setbeamercolor{block title alerted}{use=structure,fg=white,bg=titlecolor}
\setbeamercolor{section in head/foot}{bg=white,fg=shaded}
\setbeamercolor{frametitle}{fg=alerted,bg=white}
\setbeamercolor{block body}{bg=structure!10}
\setbeamercolor{block title}{bg=structure!20}
\setbeamercolor{block body example}{bg=green!10}
\setbeamercolor{block title example}{bg=green!20}

%Fonts
\setbeamerfont{subsection in toc}{size={\fontsize{7}{10}}}
\setbeamerfont{author in sidebar}{size={\fontsize{5}{8}}}
\setbeamerfont{section in sidebar}{size=\fontsize{5}{6}\selectfont}

%Theme
\setbeamertemplate{itemize item}{\scalebox{0.8}\textbullet}
\setbeamertemplate{itemize subitem}{\scalebox{0.4}\textbullet}
\setbeamertemplate{itemize items}[default]
\beamertemplatenavigationsymbolsempty
\setbeamertemplate{footline}[frame number]
\newenvironment{sitemize}{\begin{itemize}[label=$\to$]}{\end{itemize}}
\setbeamersize
{
	text margin left=1.5cm,
	text margin right=1.5cm
}

\begin{document}

	\begin{frame}
		\thispagestyle{empty}
		\begin{center}
			\textit{Applied Labor Economics}\\
			{ \huge
				\vspace{1.3cm}
				\alert{Shortage Lists and Monopsony}}\\
			\vspace{0.6cm} {\large \alert{Project Proposal}} \\
			\vspace{1.3cm}
			\textbf{Lisa Chardon--Denizot, Maria Romagnoli} \\
			\vspace{0.8cm}
		 	April 3rd, 2026
		\end{center}
	\end{frame}

\section*{Context}

\begin{frame}{Context - \textit{research question}}

\begin{columns}
\begin{column}{0.5\textwidth}

\alert{\textbf{métiers en tension}}
\begin{itemize}
    \item Easier hiring of non-EEA workers in occupations facing labor shortages.
    \item List varies by \alert{region} and \alert{occupation}. 
    \item Three updates (2008, 2021, 2025).
\end{itemize}

\vspace{0.4cm}

\alert{\textbf{hypothesis}}
\begin{itemize}
    \item Incentive for people seeking better status stability to work in shortage occupations.
    \item Employers know this $\rightarrow$ price wages down.
\end{itemize}
\end{column}

\begin{column}{0.5\textwidth}

\begin{alertblock}{research question}
Do non-EEA workers accept lower wages in shortage occupations because they give access to better legal status?
\end{alertblock}

\begin{figure}[h]
    \centering
    \includegraphics[width=\textwidth]{nanterre.png}
\end{figure}

\end{column}
\end{columns}
\end{frame}

\begin{frame}{Context - \textit{research question}}
\begin{figure}[h]
    \centering
    \includegraphics[width=0.7\textwidth]{somejobs.png}
\end{figure}
\end{frame}

\section*{Mechanism}

\begin{frame}{Mechanism - \textit{quick sketch}}

	\begin{itemize}
		\item Let's tie our reasoning to a job search model with Nash bargaining.
		\item The wage level is the result of: \\
		\vspace{0.2cm}
	
		$w = \underbrace{(1-\beta)}_{\substack{\text{firm power}}} \cdot \underbrace{\bar{w}_i}_{\substack{\text{worker's outside option}}} + \underbrace{\beta}_{\substack{\text{worker power}}}  \cdot \underbrace{p}_{\substack{\text{total surplus}}}$
	\end{itemize}

	\vspace{0.3cm}

	\begin{itemize}
	\item Three worker types: natives, EEA workers, and non-EEA workers.

	\begin{itemize}
		\item \alert{natives and eea workers:} outside option is market wage. \\ 
		\vspace{0.05cm}
		$\bar{w}_{\text{natives}} \approx \bar{w}_{\text{eea}} = \bar{w}^{*}$
		\vspace{0.2cm}
		\item \alert{non-eea workers:} outside option is undocumented work. \\
		\vspace{0.05cm}
		$\bar{w}_{non eea} = \max(\bar{w}^{\text{undocumented}}, \underbrace{w^{\text{other}}-\text{permit cost}}_{\text{changing jobs}}) \ll \bar{w}^{*}$
	\end{itemize}

	\vspace{0.2cm}
	\item Different \alert{outside options} $\rightarrow$ wage gap ? \\
	\vspace{0.2cm}
	$\Delta w \equiv w_{eea} - w_{non-eea} = (1 - \beta)(\bar{w}^* - \bar{w}_{non-eea}) > 0$
	\end{itemize}	
\end{frame}


\begin{frame}{Mechanism - \textit{different forces}}

	\alert{\textbf{Several effects at play:}}
	\vspace{0.2cm}
	\begin{enumerate}
		\item \alert{\textit{monopsony effect}}: firms post lower wages because they know non-EEA workers have a lower outside option.
		\vspace{0.1cm}
		\item \alert{\textit{discrimination effect}}: firms offer lower wages to non-EEA workers because they're discriminating $\rightarrow$ many nuanced mechanisms \textcolor{violet}{(Glover et. al, 2022)}.
		\vspace{0.1cm}
		\item \alert{\textit{shortage effect}}: firms post higher wages to attract workers in shortage occupations.
		\vspace{0.1cm}
		
	\end{enumerate}

	\vspace{0.4cm}
	\alert{\textbf{Our goal:}}
	\begin{itemize}
		\item Elicit something akin to the \alert{\textit{monopsony effect}} while dis-entangling the rest. 
	\end{itemize}
 \end{frame}



\begin{frame}{Empirics - \textit{motivation}}
    \begin{itemize}
	\item Rich empirical research suggests our mechanism could make sense. 
	\vspace{0.3cm} 
	
	\begin{itemize}
	\item \textcolor{violet}{\textbf{Signorelli (2024)}} France 2008 shortage list reform. Increase in hiring of non-EU workers. Wages fell in shortage occupations more for natives than for migrants. 
	\vspace{0.2cm}
	\item \textcolor{violet}{\textbf{Amior \& Stuhler (2023)}} Germany after Iron Curtain. Immigration strengthened firms' monopsony power at the bottom of pay distribution. Large decline in wages.
	\vspace{0.2cm}
	\item \textcolor{violet}{\textbf{Naidu, Nyarko \& Wang (2016)}} United Arab Emirates. Firms paid workers less than their marginal product. Reform untied workers from employers, wages rose.
	\end{itemize}

	\vspace{0.3cm}
	\item \alert{Our contribution:} use more recent data than Signorelli (2024), try to explicitely isolate the \alert{\textit{monopsony effect}} from taste-based discrimination. 
\end{itemize}
\end{frame}

\section*{Empirics}

\begin{frame}{Empirics - \textit{data}}

\alert{\textbf{ForCE}}   
   \begin{enumerate}
    \item FH table $\rightarrow$ job-seeker records from \textit{France Travail}.
        \begin{itemize}
            \item socio-demographics.
            \item reservation wage.
        \end{itemize}
    \item MMO table $\rightarrow$ private sector work contracts from \textit{Déclaration Sociale Nominative}.
    \begin{itemize}
            \item wage.
            \item occupation code.
        \end{itemize}
   \end{enumerate}

\vspace{0.2cm}

\alert{\textbf{Shortage Lists}} 
 \begin{enumerate}
    \item Digitalized \textit{arrêtés} with FAP and ROME. 
   \end{enumerate}
\end{frame}


\begin{frame}{Empirics - \textit{strategy}}

\alert{\textbf{We need to:}}
\begin{enumerate}
	\item Retrieve the causal effect on wages of being non-EEA in a listed occupation. 
	\item Make sure we're not capturing capturing taste-based discrimination. 
\end{enumerate}
\vspace{0.4cm}
\alert{\textbf{We plan on doing this by:}}
\begin{enumerate}
	\item A nice DiD specification.
	\item Using only certain ethnic groups.
\end{enumerate}
\end{frame}

\begin{frame}{Empirics - \textit{strategy}}
\label{DiD}

\alert{\textbf{Baseline DiD}}
\begin{align*}
\ln(w_{ijrt}) &= \beta_0 + \beta_1 \text{NonEEA}_{i} + \beta_2 \text{EEA}_{i} + \beta_3 \text{Listed}_{jrt} 
+ \beta_4 (\text{NonEEA}_{i} \times \text{Listed}_{jrt}) \\ &+ \beta_5 (\text{EEA}_{i} \times \text{Listed}_{jrt} 
+ \mathbf{X}_{ijrt}'\boldsymbol{\gamma} + \text{fixed effects (tbd)}+ \varepsilon_{ijrt}
\end{align*}

\begin{itemize}
	\item \alert{explained var}: $w_{i,j,r,t}$, wage of worker $i$ in occupation $j$, region $r$, at time $t$.
	\item \alert{fixed effects}: either \textit{occupation $\times$ region $\times$ year} (but then no $\beta_3$), or a combination.
	\item \alert{controls}: age, gender, experience, tenure, full-time. 
\end{itemize}
\vspace{0.3cm}

\alert{\textbf{Key parameters}}
\vspace{0.1cm}
\begin{itemize}
	\item omitted group: natives in unlisted occupations.
	\item \alert{monopsony} $\rightarrow$ $\beta_4 - \beta_5$, differential listing effect on wages for non-EEA vs. EEA (compared to natives).
\end{itemize}
\hyperlink{ethnic}{\beamerbutton{comparing NATION groups}}
\end{frame}



\begin{frame}{Empirics - \textit{assumptions and limitations}}
    \begin{itemize}

    \item \alert{\textbf{Parallel trends}} \\
		In the absence of listing, wage trends for non-EEA workers in eventually-listed occupations would have been parallel to natives/EEA. \\
		$\Rightarrow$ \textit{bias could go both ways, depending on pre-trends.}
		\vspace{0.1cm}
	
	\item \alert{\textbf{Spillovers}} \\
		Listing in one occupation does not affect wages in other unlisted occupations. \\
		$\Rightarrow$ \textit{bias could go both ways, depending on spillover sign.}
		\vspace{0.1cm}

	\item \alert{\textbf{Composition}} \\
		This policy does not change the type of non-EEA workers entering the occupations. \\
		$\Rightarrow$ \textit{assuming positive selection, bias downwards.}
		\vspace{0.1cm}

	\item \alert{\textbf{Shock exogeneity}} \\
		Listing is exogenous.  \\
		$\Rightarrow$ \textit{it is not, might conflate shortage premium with monopsony markdown. DiD solves this, if no heterogeneity in shortage premia.}
		\vspace{0.1cm}
    \end{itemize}

\end{frame}

\section*{Conclusions}
\begin{frame}{Conclusions - \textit{our concerns}}

\alert{\textbf{Time}}
\begin{enumerate}
	\item Not managing to clean/wrangle the data appropriately !
\end{enumerate}
\vspace{0.2cm}

\alert{\textbf{Design}}
\begin{enumerate}
	\item Do we need to "purge" the wages from time and location effects ?
	\item Do we have an appropriate design given the variation ? Listed occupations change over time and over region, but not many of them.
\end{enumerate}
\vspace{0.2cm}


\alert{\textbf{Econometrics}}
\begin{enumerate}
	\item Enough statistical power with the reduced samples ?
	\item If we choose a mix of FE, what variation do we end up using ? 
\end{enumerate}
\end{frame}

\begin{frame}
		\thispagestyle{empty}
		\begin{center}
			\textit{}\\
			{ \huge
				\vspace{1.3cm}
				\alert{Thank you for listening!}}\\
			\vspace{0.6cm} {\large \alert{Any questions or feedback ?}} \\
			\vspace{1.3cm}
			\textbf{} \\
			\vspace{0.8cm}
		\end{center}
	\end{frame}


\section*{Appendix}
\begin{frame}{Appendix - \textit{focusing on monopsony}}
\label{ethnic}

\begin{itemize}
    \item How to separate the monopsony channel from taste-based discrimination ?
    \item If possible, compare groups similar on ethnic, cultural, religious lines.
\end{itemize}
\vspace{0.4cm}

{\centering
\renewcommand{\arraystretch}{1.2}
\begin{tabular}{lll}
\toprule
\textbf{Non-EEA (Count)} & \textbf{EEA (Count)} & \textbf{Bias?} \\
\midrule
Yugoslavs (34,000) & Croats + Slovenes (3,873) & Partial \\
Yugoslavs (34,000) & Bulgarians + Romanians (164,859) & Partial \\
Turks (111,043) & RO + BG + GR + CY (171,348) & Religious \\
Other Europeans (133,339) & Baltics + Central EU (61,839) & None \\
\bottomrule
\end{tabular}
}

\vspace{0.3cm}
{\small 
\textbf{Note:} Groups in the final row are Slavic (except Estonia and Hungary) and Christian. Need to assume no discrimination against Orthodoxy.
}  

\hyperlink{DiD}{\beamerreturnbutton{Back}}
\end{frame}

\end{document}
